\section{Introduction}
\label{sec:intro}

Markdown ``skills'' are a real distribution channel. The Agent Skills
standard fixes a \texttt{SKILL.md} format and a shared
\texttt{.agents/skills/} discovery path, and its own client showcase listed
forty-six adopting products in August 2026, among them Claude Code, Codex,
Cursor, GitHub Copilot, VS Code, Gemini CLI and
Amp~\citep{agentskills2025}. These files are installed and running inside
coding agents today. Almost none of them carry evidence that they help.

A second layer has grown on top of the first. Engines now rewrite skill
files automatically: GEPA proposes revised instructions from reflection over
rollouts~\citep{gepa2026}, and Microsoft's SkillOpt makes bounded edits to a
skill body and keeps any edit that scores strictly
better~\citep{skillopt2026}. Both accept an edit on a single unreplicated
score comparison, with no notion of a difference that could have arisen by
chance.

This paper is about the two controls that acceptance rule is missing, and
about what happens when they are added. The first is a placebo: no engine in
this space compares its output against a document of matched length and
structure whose content says nothing, so a gain from an evolved skill and a
gain from any structured text in the system prompt are indistinguishable in
every published number. The second is power arithmetic computed from the
design before the run: how small a $p$-value the design can produce at all,
and how large an effect it could detect. Both are cheap, and this study is
evidence for both: we ran it with the placebo in place and computed the
arithmetic only afterwards.

\paragraph{What we did.}
We held the model fixed, evolved a decision-making skill with each engine
against a training split, froze both winners by content hash, and then
minted a fresh holdout the engines had never seen. Seven arms ran across two
deterministic passes over that holdout (14,700 calls in all): no skill
(\texttt{off}), the initial seed skill (\texttt{on}), a placebo matched to
that seed skill on word count and structure (\texttt{placebo}), length- and
structure-matched placebos for each evolved winner (\texttt{placebo-gepa} and
\texttt{placebo-skillopt}), and the two evolved winners (\texttt{gepa} and
\texttt{skillopt}), closing the length confound (\Cref{sec:length-confound}).
Every comparison is McNemar's exact test against each arm's matched placebo,
paired by item, with Holm correction across a family of \NUM{\familySize} per
item set, on \NUM{\unseenItems} unseen items across \NUM{\templatesUnseen}
held-out templates and \NUM{\seenItems} seen items across \NUM{\templatesSeen}
trained templates, analysed separately, backed by template-unit cluster
sign-flip permutation tests. The prediction was committed to git before any data
existed, and the run's provenance gate refuses a result whose prediction cannot
be shown by ancestry to predate it.

\paragraph{What we found.}
No arm rejects, on either item set, after correction, at the pre-registered
dual bar (Holm $q < 0.05$ and template cluster sign-flip $p < 0.0167$).
At the item level, both evolved winners appear to score decisive victories on
one set: SkillOpt reaches $\NUM{\effectUnseenSkilloptVsPlaceboSkillopt}$ against
its matched placebo on unseen scenarios (raw $p = \NUM{\pUnseenSkilloptVsPlaceboSkillopt}$,
Holm $q = \NUM{\qUnseenSkilloptVsPlaceboSkillopt}$), and GEPA reaches
$\NUM{\effectSeenGepaVsPlaceboGepa}$ on seen scenarios (raw $p = \NUM{\pSeenGepaVsPlaceboGepa}$,
Holm $q = \NUM{\qSeenGepaVsPlaceboGepa}$). But under template-unit clustering,
neither survives: SkillOpt yields $p = \NUM{\clusteredUnseenSkilloptVsPlaceboSkillopt} > 0.0167$
and GEPA yields $p = \NUM{\clusteredSeenGepaVsPlaceboGepa} > 0.0167$.
Their gains are concentrated in specific template structures rather than uniform
across tasks.

The initial seed skill falls at or below the placebo on both sets
(\NUM{\effectUnseenOn} unseen, \NUM{\effectSeenOn} seen). Both winners had
written training-item constants into their own bodies, each asserting specific
thresholds and worked examples lifted verbatim from training items
(\Cref{sec:whatengines}). Prompting text provoked runaway generations that ran
to the output cap, causing the empty prompt (\texttt{off}) to outperform the
seed skill (\texttt{on}) by $\sim 6$--$7$ points across both sets.
Control tokens (\texttt{/think} or \texttt{/no\_think}) appeared on only 48 of
6,860 calls (0.70\%) and all 48 parsed cleanly (\Cref{sec:parsefail}).
Determinism was complete: pass-to-pass agreement was 100\% across all seven arms,
and the A/A control returned \NUM{\aaPairs} of \NUM{\aaPairs} identical.

\paragraph{Contributions.}
\begin{enumerate}
  \item \textbf{The first placebo-controlled evaluation of automated skill
    optimisation we are aware of}, with dedicated length- and structure-matched
    placebos for each evolved winner and an A/A determinism control reported
    beside it (\Cref{sec:results}). The placebos match or beat the skills on
    both item sets.
  \item \textbf{Template-level permutation clustering that eliminates false
    discoveries.} Expanding the unseen set to \NUM{\templatesUnseen} templates
    lowered the theoretical permutation floor to $2^{-7} = 0.0078$, giving the
    design power to detect genuine effects. The dual bar demonstrates that
    item-unit tests would have falsely claimed significance for both engines,
    while cluster sign-flip tests correctly identify that the gains are
    template-specific (\Cref{sec:clustered}).
  \item \textbf{A discrimination measure reported beside accuracy}, showing
    that only \NUM{\lowSignalTemplates} of \NUM{\studyTemplates} templates has
    low informedness ($J < \NUM{\lowSignalThreshold}$) across arms, and that
    \NUM{\signalItems} items carry clean discriminative signal (\Cref{sec:signal}).
  \item \textbf{Evidence of systematic memorisation across optimisers.} Both
    engines wrote training constants and cutoffs into their skill bodies,
    demonstrating that score ratchets without cluster-aware validation readily
    exploit training artifacts (\Cref{sec:whatengines}).
\end{enumerate}

\paragraph{Scope, stated once and up front.}
This is one target model, one corpus, one harness. The target is a 1.7B
parameter local model, and every larger model we later screened solves the
corpus with an empty prompt, so no replication above that tier exists
(\Cref{sec:noreplication}). The claim is that on this model, under these
controls, neither engine's output beat a matched placebo. Whether it would on a
frontier model is open, and the screen that says why is itself a result:
every hosted model we measured solves this corpus with an empty prompt
(\Cref{sec:noreplication}).
